\begin{frame}{Schéma de l'index }


 Le schéma donne la liste de tous les champs d'un document 

\vfill

Nombreuses options : 
      \begin{itemize}
        \item type (numérique, entier)
        \item possibilité de calcul de la valeur du champ à partir d'un autre
        \item traitements divers sur les valeurs du champ
      \end{itemize}

\end{frame} % ending frame Schéma de l'index 

\begin{frame}[fragile]{Reprenons un  document bien connu}
  \begin{minted}[baselinestretch=.9, fontsize=\footnotesize]{json}
{
  "title": "Vertigo",
  "year": 1958,
  "genre": "drama",
  "summary": "Scottie Ferguson, ancien inspecteur de police... ",
  "director": 	{
    "_id": "artist:3",
    "last_name": "Hitchcock",
    "first_name": "Alfred",
    "birth_date": "1899"	
  },
  "actors": [
     {
       "_id": "artist:15",
       "first_name": "James",
       "last_name": "Stewart",
       "birth_date": "1908",
       "role": "John Ferguson"	
     }
  ]
}
\end{minted} 
\end{frame}

\begin{frame}[fragile]{Son schéma}
	\begin{minted}[baselinestretch=1.0, fontsize=\footnotesize]{json}
{ "mappings": {
    "movie": {
      "_all":       { "enabled": true  },
      "properties": {
        "title":    { "type": "string"  },
        "year":     { "type":   "date", "format": "yyyy" },
        "genre":      { "type": "string" },
        "summary": { "type": "string"},
        "director":     {
          "properties": {
            "_id": { "type":"string" },
            "last_name": { "type":"string"},
            "first_name": { "type":"string"},
            "birth_date": { "type": "date", "format":"yyyy"}
          }},
        "actors": { 
          "type": "nested",
          "properties": {"..."}
        }
      }
    }
  }
}
  \end{minted}
\end{frame}

\begin{comment}
\begin{frame}[fragile]{Détaillons le schéma}
  \begin{itemize}
    \item la liste des champs, dans l'élement ``mappings`` 
    \item le champ \verb+_all+ est spécial, tout y est par défaut concaténé, de
      façon à permettre une recherche sur tous les champs
    \item la liste des {\bf types de champ}, dans l'élément ``properties``.
  \end{itemize}
\end{frame}
\end{comment}

\begin{frame}{Définition des types et de la clé}

Chaque champ utilisé dans un document doit apparaître dans un
des éléments. Sinon, ElasticSearch propose un
      "Dynamic Mapping"
      
      \vfill
   ElasticSearch fournit tout un ensemble de types pré-définis pour      
      les besoins courants; 

\vfill

Les options indiquent d'éventuels traitements à appliquer                      
 à chaque valeur du type avant son insertion dans l'index
 
 \vfill
 
 \begin{blocgauche}
Ex: type \variable{string} (remplacé par \variable{text} depuis 5.0)
      \begin{itemize}
        \item on lui définit un ``analyzer''  ``standard''
        \item traite le texte pour une recherche
          plein-texte
        \item retenir :  permet d'indexer chaque  mot, et donc de faire
          des recherches sur toutes les combinaisons de mots
      \end{itemize}
\end{blocgauche}

\end{frame}



\begin{frame}[fragile]{Définition des champs}
Le {\bf nom} et le {\bf type} sont les informations de base
\begin{minted}[frame=lines, baselinestretch=.9, fontsize=\footnotesize, framesep=2mm]{json} 
{
  "genre": { 
		"type": "string",
		"index": "no",
		"store": "true",
		"index_options": "offsets"
     }
}
\end{minted}

\vfill
Ensuite, divers attributs (souvent optionnels) :
    \begin{itemize}
      \item \variable{index} indique simplement si le champ peut être utilisé dans une recherche;  
      \item \variable{store} indique si la \alert{valeur} du champ est \alert{stockée} dans l'index,
         
\begin{comment}
          \item remarquez la structure pour le réalisateur (\variable{directors})
          \item enfin, \variable{nested} pour les champs ayant plusieurs
       
    valeurs, soit, concrètement, un {\bf tableau} en JSON; c'est le
    cas par exemple pour le nom des acteurs.
\end{comment}
    \end{itemize}
\vfill
   
\begin{blocgauche}
 L'option \variable{store} permet de considérer ElasticSearch comme une base de données
 indexées.  
\end{blocgauche}
\end{frame}

\begin{frame}[fragile]{Définition des champs}
 Les attributs \variable{index}  et \variable{store} sont très importants
\vfill
    \begin{itemize}
      \item \variable{index=true}, \variable{store=false}: on pourra interroger
      le champ, mais il faudra accéder au document principal dans la  base documentaire si on veut sa valeur;
    \item \variable{index=true}, \variable{store=true}: on pourra interroger
      le champ, \alert{et} accéder à sa valeur dans l'index;
    \item \variable{index=false}, \variable{store=true}: on ne peut pas interroger le champ, mais
      on peut récupérer sa valeur dans l'index.
    \end{itemize}

\begin{blocgauche}
\variable{index=false}, \variable{store=false}: n'a pas de sens à priori.
\end{blocgauche}
\end{frame}




\begin{frame}[fragile]{Définition des champs (suite)}
Comment peut-on indexer un champ sans le stocker ?
  \begin{itemize}
    \item c'est notamment le cas pour les textes qui sont décomposés en
      {\bf termes}: chaque terme est indexé indépendamment
    \item très difficile pour l'index de reconstituer le texte
  \end{itemize}

\vfill

Doit-on stocker les documents dans l'index? Question de compromis: 
\begin{itemize}
  \item (-) {\bf stocker} une valeur prend plus d'espace que {\bf l'indexer}; 
  À l'extrême,  la base documentaire
    est dupliquée dans l'index
  \item (+) Peut considérablement simplifier l'architecture
\end{itemize}

\vfill

\begin{blocgauche}
Attention aux autres facteurs, comme la latence pour les insertions
et mises à jour.
\end{blocgauche}

\end{frame}

\begin{frame}[fragile]{Champ par défaut, document original}

Le champ \variable{\_source} contient le document, mais n'est pas indexé. Il est
possible ainsi de récupérer la totalité du document passé à l'indexation. 

\vfill

Le champ \variable{\_all} contient la concaténation de tous les autres champs; il
est indexé, pas stocké.
      
\vfill
\begin{blocgauche}
Permet de
retrouver les documents contenant une valeur, sans savoir dans quel champ elle
se trouve. 
\end{blocgauche}
\end{frame}

\begin{comment}
\begin{frame}[fragile]{Champ calculé}
 Il est possible de créer des champs \verb+all+ spécifiques, correspondants à
 des besoins précis, à l'aide d'une instruction \verb+copy_to+
      
\vfill
      
  \begin{itemize}
    \item Ainsi, le nom et le prénom peuvent être regroupés au sein d'un champ ``nom
complet``. 
    \item Cela permet aussi d'indexer certains champs plusieurs fois, de différentes
façons.
    \item par exemple le contenu d'un titre
    est indexé comme une chaîne de caractères dans le champ \variable{title}, et comme un texte
    "tokenisé" quand on le copie dans le champ \variable{text};
  \end{itemize}
\end{frame}




\begin{frame}[fragile]{Exemple de champ calculé}

\begin{minted}[ baselinestretch=1.2, fontsize=\footnotesize,]{json}
{
  "mappings": {
    "my_type": {
      "properties": {
        "first_name": {
          "type": "text",
          "copy_to": "full_name" 
        },
        "last_name": {
          "type": "text",
          "copy_to": "full_name" 
        },
        "full_name": {
          "type": "text"
        }
      }
    }
  }
}
\end{minted}
\end{frame}

\end{comment}


\begin{frame}[fragile]{Configuration des pré-traitements}

Analyseur = une chaîne de traitement constituée de tokeniseurs et de filtres

\begin{minted}[fontsize=\footnotesize]{bash}
curl -XPUT '<url>/monindex' -H 'Content-Type: application/json' -d @settings.json
\end{minted}
\vfill

Contenu de \texttt{settings.json}:

\begin{minted}[fontsize=\footnotesize]{json}
{"settings": {
    "analysis": {
      "analyzer": {
        "custom_lowercase_stemmed": {
          "tokenizer": "standard",
          "filter": [
            "lowercase",
            "custom_english_stemmer"
          ]
        } 
     } 
   } 
}
\end{minted}

\end{frame}

\begin{frame}[fragile]{Test d'un analyseur}

On peut tester un analyseur en lui soumettant un texte-test.

\begin{minted}[fontsize=\footnotesize]{bash}
curl -XPUT '<url>/monindex/_analyze' -H 'Content-Type: application/json' -d @test.json
\end{minted}
\vfill

Contenu de \texttt{test.json}:

\begin{minted}[fontsize=\footnotesize]{json}
{
        "analyzer": "custom_lowercase_stemmed",
        "text":"Finiras-tu ces analyses demain ?"
}
\end{minted}

\end{frame}

\begin{frame}[fragile]{Résultat d'un test}


\begin{minted}[fontsize=\footnotesize]{json}
{
  "tokens" : [ {
    "token" : "finira",
    "start_offset" : 0,
    "end_offset" : 7,
    "type" : "<ALPHANUM>",
    "position" : 0
  }, {
    "token" : "tu",
    "start_offset" : 8,
    "end_offset" : 10,
    "type" : "<ALPHANUM>",
    "position" : 1
  }, 
  "...": "..."
}
\end{minted}

\end{frame}


\begin{frame}{Résumé}

La configuration d'un moteur de recherche, c'est:

\begin{itemize}
 \item Définir le schéma des champs indexés (stockage / indexation / les deux)
 \item Définir les chaînes de traitement.
\end{itemize}

\vfill

C'est \alert{aussi} beaucoup d'autres choses.

\begin{itemize}
 \item Définir les synonymes
 \item Définir les mots protégés (acronymes)
 \item Définir les mots vides
 \item Acquérir et exploiter les actions utilisateurs
\end{itemize}


\vfill
\begin{blocgauche}
À pratiquer et approfondir (projet). 
\end{blocgauche}

\end{frame}

\end{document}

\begin{comment}
\begin{frame}[fragile]{Schéma}
On peut stocker le schéma dans un fichier json et créer l'index dans
ElasticSearch comme ceci : 
\begin{minted}[ baselinestretch=1.2, fontsize=\footnotesize,]{bash}
    curl -XPUT '<url>/monindex' -H 'Content-Type: application/json' -d movie-schema.json
  \end{minted}
\end{frame}


\begin{frame}[fragile]{Mise à jour de l'index}
\begin{itemize}
  \item Avec ElasticSearch, il est (plus) difficile de faire évoluer un schéma (qu'avec BDD classique)
\item Sauf rares exceptions, il faut en général créer un nouvel index et réindexer les données
\item ElasticSearch permet la copie d'un index vers l'autre avec l'API \verb+reindex+
\end{itemize}

\begin{minted}[frame=lines,
               baselinestretch=.9,
               fontsize=\footnotesize,
               framesep=2mm]{bash}
 curl -XPOST 'localhost:9200/_reindex?pretty' -H 'Content-Type: application/json' -d'
{
  "source": {
    "index": "nfe204"
  },
  "dest": {
    "index": "new_nfe204"
  }
}'

\end{minted}
\end{frame}

\end{comment}
